% SPDX-License-Identifier: Apache-2.0
\section{Signals, Ports and Interfaces}
\label{sec:e-wiring}

These three were the muddiest part of the merged specification, and LHDL's
version of them is where its grammar and its examples disagreed most. The
embedding settles them, and the settlement is one idea.

\begin{rulebox}[Direction is not a property of a wire]
The same wire is driven at one end and read at the other. Direction
belongs to the end, not to the wire, so a unit never holds a wire. It
holds an end.
\end{rulebox}

\subsection{A signal has two ends}

Creating a signal yields its ends separately, the way
\code{std::sync::mpsc} yields a sender and a receiver.

\begin{lstlisting}[caption={A signal is created as a pair. \code{T} is a struct, so a wire may be a bundle without a second concept for the bundled case.},label={lst:split}]
pub fn signal<T: Copy + Default>() -> (Out<T>, In<T>);

let (tx, rx) = signal::<Beat>();
\end{lstlisting}

\code{Out<T>} has \code{set} and no \code{get}. \code{In<T>} has \code{get}
and no \code{set}. A unit given the reading end cannot drive, and the
compiler says so in the plainest possible terms:
\code{error[E0599]: no method named `set` found for struct `In`}.

\subsection{Cardinality is the opposite of mpsc, and the types say so}

A channel in software has many senders and one receiver. A wire has one
driver and many readers. The library inverts the derive accordingly.

\begin{table}[H]
\caption{Why \code{In} is \code{Clone} and \code{Out} is not.}
\label{tab:cardinality}
\centering
\footnotesize
\begin{tabular}{@{}lll@{}}
\toprule
 & Sender or driver & Receiver or reader \\
\midrule
\code{mpsc} & \code{Clone}     & unique \\
a wire      & unique           & \code{Clone} \\
\bottomrule
\end{tabular}
\end{table}

That one line of derive buys the oldest error message in hardware design.
A second driver on a net would be a second \code{Out}, \code{signal}
returns exactly one, and \code{Out} is not \code{Clone}, so there is no
way to make another.

\begin{lstlisting}[caption={Two drivers on one wire. The classic multiple-driver error is a move error.},label={lst:twodrivers}]
let (tx, _rx) = signal::<u32>();
let a = Driver { out: tx };
let b = Driver { out: tx };
// error[E0382]: use of moved value: `tx`
\end{lstlisting}

Nothing in the language had to define that rule. Fanout, meanwhile, is
\code{rx.clone()} and costs nothing, because reading a wire is free.

\subsection{An interface is declared, not written out}

An interface is a named group. Its members are signals and channels, and
it holds no direction, because its members do not. A \term{role} pairs
each member with the end this side takes, and a \term{port} is what a unit
is handed.

Written by hand, that is one struct for the interface and one more per
role, with every member repeated in each. It is mechanical, and it is also
where an interface and its roles drift apart. So it is declared.

\begin{lstlisting}[caption={Complete interface declaration and its two roles; all supporting definitions are generated automatically.},label={lst:iface}]
interface! {
    Wishbone {
        adr: Signal<U<32>>,
        ack: Signal<Bit>,
        dat: Chan<Beat>,
    }
    role Initiator { out adr, in ack, out dat }
    role Target    { in adr,  out ack, in dat }
    role Monitor   { in adr,  in ack,  in dat }
}
\end{lstlisting}

That expands to the interface, a port struct per role with each member's
end already the right type, and a constructor yielding one port for each.
\code{interface!} is a procedural macro, and the reason it is not a
\code{macro\_rules!} one is the next subsection.

\begin{lstlisting}[caption={Using it. Driving the wrong end is a missing method, not a convention.},label={lst:iface-use}]
let (m, t, mon) = Wishbone::new();

m.adr.set(U::new(0x1000)); // Initiator drives adr
t.ack.set(Bit::One);       // Target drives ack
let a = t.adr.get();       // and reads it at the other end
let w = mon.adr.get();     // as does the monitor
\end{lstlisting}

The macro never learns what kind a member is. Every member implements one
trait, \code{Member}, with \code{Driver} and \code{Reader} associated
types and a \code{split} that consumes the member and returns both ends.
A \code{Signal} splits into \code{Out} and \code{In}; a \code{Chan} splits
into \code{Tx} and \code{Rx}. The macro writes
\code{<Ty as Member>::Driver} and lets the trait decide, so adding a third
kind of member changes no macro.

Constructing the interface yields every port at once, which is what keeps
the guarantee of Section~\ref{sec:e-wiring}. In the constructor a reader
is cloned, because readers fan out, and a driver is moved. Three roles
may all read \code{adr}; two may not both drive it.

\subsubsection{Why it is a procedural macro}

It was a \code{macro\_rules!} macro first, and that version compiled
with three defects that a declaration should not be able to have. It
spelled the direction \code{inp}, because \code{in} is a keyword and a
pattern cannot match one as a word. It fixed the number of roles at two.
And a role that forgot a member compiled, silently, into a port without
that member.

A procedural macro reads the declaration as data rather than matching it
as a pattern, and all three go. \code{in} is an identifier to it. Roles
are a list. And it checks, before emitting anything, that every role
names every member exactly once and that no member is driven twice:

\begin{lstlisting}[style=ascii,caption={What the declaration cannot be. Both are the macro's own messages; the second would otherwise arrive as a move error naming a temporary.},label={lst:iface-errors}]
error: role `Target` does not say which end of `dat` it takes
error: member `adr` is driven by more than one role: A, B
\end{lstlisting}

The macro is under two hundred lines and uses neither \code{syn} nor
\code{quote}, because the grammar is four productions. The
\code{macro\_rules!} version is kept as probe~18, as the record of where
that route stopped, and Section~\ref{sec:e-walls} gives the other place
it stops. Both point the same way as Section~\ref{sec:e-elaboration}.

\subsection{The clock is in the type}

Following RHDL~\cite{rhdl}, a signal names the clock domain it is in as a
second type parameter.

\begin{lstlisting}[caption={A signal knows its clock. \code{Clock} is a trait with a name and, from Section~\ref{sec:e-phase}, a period and a phase; the default is a named clock rather than an absence of one.},label={lst:clock}]
pub trait Clock: 'static {
    const NAME: &'static str;
    const PERIOD: u64 = 2;   // time steps between edges
    const PHASE: u64 = 0;    // time step of the first edge
}

pub struct DefaultClock;
impl Clock for DefaultClock { const NAME: &'static str = "clk"; }

pub struct Out<T, C: Clock = DefaultClock>(..);
pub struct In<T, C: Clock = DefaultClock>(..);
pub struct Reg<T, C: Clock = DefaultClock>(..);
\end{lstlisting}

The default is the decision worth explaining. It is not \code{()}, and
it is not ``no clock''. It is \emph{the} clock: the one a single-clock
design has, which is most designs. A single-clock design then never
mentions a domain, \code{In<u32>} is a complete type, and a two-clock
design names exactly the second clock. \code{()} would say nothing is
there, and a register cannot live in nothing.

The default is also a domain like any other. It does not unify with
\code{Clk400}, so a signal from the default clock handed to a port in
another is refused, and so is any other pairing:

\begin{lstlisting}[style=ascii,caption={A domain crossing without a crossing. Both are refused, the second showing the default is not a wildcard.},label={lst:clock-err}]
expected `In<u32, Clk400>`, found `In<u32, Clk100>`
expected `In<u32, Clk400>`, found `In<u32>`
\end{lstlisting}

Values cross clock domains through explicit channel bridges such as
\code{ChanCdc} in \code{//lib/parts}, which accepts transactions on one
clock and delivers them on the other (issue 1017). Internally,
\code{ChanCdc} implements an asynchronous dual-clock FIFO: write and
read pointers are Gray-coded and synchronized across domains via
two-stage flip-flop synchronizers. Raw wires cannot bridge domains
directly because signal endpoints must belong to matching clock types,
preventing unmitigated domain crossings. While the merged specification
relied on compiler tag analysis to flag clock domain mismatches, the
embedded runtime enforces this invariant through the Rust type system.

A unit that does not care which clock it is in says so by being generic
over \code{C: Clock}, which is the Rust way of saying it and costs the
unit nothing.

\subsubsection{Clocks and tags are different things}

Section~\ref{sec:e-binding} gave a \code{Tag}: a synchronisation domain
with a policy, handshake and capacity. That is LHDL's idea, and it is
about how data moves. A \code{Clock} defines the physical switching edge
that triggers register updates. A single clock domain can host multiple
synchronisation tags, accommodating elastic pipelines and raw synchronous
paths concurrently. The system configuration maps tags onto physical
clocks. Decoupling these concepts avoids forcing every elastic domain to
instantiate a distinct hardware clock network.

\subsubsection{Two clocks stand in a stated relation}
\label{sec:e-phase}

A design with two clocks depends on how they stand to each other, and
on nothing else about them. The frequency in hertz is a physical fact
and belongs to the configuration, as Section~\ref{sec:e-main} put it
there. The ratio of the two periods and the offset between their edges
are what a FIFO between them is built against, and those are stated on
the clock: a period and a phase, in ticks, a unit common to every
clock of the design. A clock with period 4 and phase 2 rises at 2, 6,
10 and falls at 4, 8, 12: the falling edge is half a period after the
rising one, which is why the default clock has period 2 rather than 1,
so that its falling edge is a tick of its own. The defaults are 2 and
0, so a single-clock design mentions neither. This is the same
statement SDC makes with \code{create\_clock -period -waveform}, and it
was chosen for that reason: it is the form a constraint file already
takes.

A register is then in a clock, \code{Reg<T, C>}, with the same default
as a wire, and a process in that clock waits for \code{C::rising()} or
\code{C::falling()}. The executor runs time in ticks and each clock has
its edges at the ticks its period and phase name; the waits built on
an edge, a channel's \code{wait} on the rising edge and \code{until} on
whichever edge it is given, are in the channel's or the named clock,
Section~\ref{sec:e-regread}. An operator's cycle is a cycle of whatever
clock the process last crossed an edge of, which is the only sense a
latency has in a process.

One thing the executor does not do is check that a process stays in
its domain. A process in one clock that reads a register of another
other than through \code{ChanCdc} is a hazard, and the prototype runs
it rather than refusing it, because the register's type says which
clock it is in and the process's type says nothing. That is also the
hazard \code{ChanCdc} commits on purpose: its memory is written by one
clock and read by the other, and the pointer discipline is what makes
that safe. It is why the runtime has a \code{Mem<T, N, C>} as a type
rather than an array of registers: a memory's read port is a wire, and
the discipline is the designer's. \code{ex\_fifo.rs} in the examples
document~\cite{examples} runs it between a producer clock of period 4
and a consumer clock of period 6 and phase 2; its trace shows the
producer with nothing to push at the edge the FIFO is full, which is
the backpressure crossing the domains the other way.

\subsection{Where the ends go}

A unit does not store a reference to the interface. It is handed its port
when \code{run} is called, which is what the inputs and outputs of
\code{Unit} are for, and what makes wiring a call rather than a graph of
back-references.

\begin{lstlisting}[caption={The parent creates the interface and hands each child the end it takes.},label={lst:wire-up}]
impl Unit<(), ()> for Top {
    async fn run(&mut self, _i: (), _o: ()) {
        let (m, t) = Wishbone::new();
        join2(
            self.cpu.run(m, ()),
            self.ram.run(t, ()),
        ).await;
    }
}
\end{lstlisting}

\subsubsection{An unregistered channel is a hop that costs no cycle}

A channel between two units is registered on both sides, so a
transaction takes a cycle to cross it, and a request that crosses ten
of them on its way to a memory and back waits ten cycles for the
crossings alone. Marking a channel \code{\#[unregistered]} where a
unit of units makes it takes the register off its forward side: while
the channel is empty, the receiver sees the sender's \code{valid} and
\code{data} in the cycle they are offered, and \code{ready} stays a
register (issue 1293). Listing~\ref{lst:wire-unreg} is the chain of
\code{ex\_unreg}, whose two inner channels are unregistered; the same
chain registered answers two cycles later.

\lstinputlisting[rangebeginprefix=//\ begin\{,rangebeginsuffix=\},%
rangeendprefix=//\ end\{,rangeendsuffix=\},includerangemarker=false,%
linerange=chains-chains,caption={Two channels a unit
of units makes, marked unregistered, each sender joined before its
receiver: from \code{lib/examples/ex\_unreg.rs}.},label={lst:wire-unreg}]{lib/examples/ex_unreg.rs}

An unregistered hop is not free in time. The receiver's logic now
follows the sender's in one cycle, so every one lengthens a path, and
it is for hops whose two ends are short: a stage that only forwards,
or one that answers from a register. Two rules keep the run and the
netlist in step. The receiver has to run after its sender in the join,
since the run polls each unit once a step, and a run that joins them
the other way stops naming the channel. A ring of wires through the
children closed by an unregistered channel has no such order, and it
is refused when the unit is lowered, naming the ring.

\begin{table*}[t]
\caption{The four concepts, and what each one is responsible for. Reading
down the table is the life of a wire.}
\label{tab:wiring}
\centering
\footnotesize
\begin{tabular}{@{}llll@{}}
\toprule
Concept & Is & Knows about direction & Held by \\
\midrule
\code{Signal<T>}, \code{Chan<T>} & one wire, or one handshaked channel & no & nobody \\
\code{Out<T>} / \code{In<T>}, \code{Tx<T>} / \code{Rx<T>} & one end of one member & yes, by which methods exist & a unit \\
interface & a named group of members & no & the unit containing both ends \\
role and port & a group of ends, one per member & yes, by construction & passed to \code{run} \\
\bottomrule
\end{tabular}
\end{table*}
